\documentclass[../main.tex]{subfiles}

\begin{document}
\begin{CJK*}{UTF8}{gbsn}

\section*{Exercise 6}

For $z \in \mathbb{R}^n$,
write $\hat{F}_z(v) = \frac{1}{n} \sum_{i=1}^n \chi_{\{z_i \leq v\}}$ 
for the empirical cumulative distribution function of $z$
with $\hat{F}_z(-\infty) = 0$ and $\hat{F}_z(+\infty) = 1$.
Write 

\begin{equation*}
  \text{Quantile}(z; \tau) 
  = \inf \curly{v \in \mathbb{R} \mid \hat{F}_z(v) \geqslant \tau}
\end{equation*}
for the empirical $\tau$-quantile of $z$ for any $\tau \in [0, 1]$.
Consequently, $\text{Quantile}(z; 0) = -\infty$.

Prove that, for any $z \in \mathbb{R}^n$,
the following statements are true:

\begin{enumerate}
  \item $\hat{F}_z(v) = \sup 
  \curly{\tau \in [0, 1] \mid \text{Quantile}(z; \tau) \leqslant v}$ for all $v \in \mathbb{R}$.
  \item $\text{Quantile}(z; \hat{F}_z(v)) \leqslant v$ for all $v \in \mathbb{R}$.
  \item $\hat{F}_z(\text{Quantile}(z; \tau)) \geqslant \tau$ for all $\tau \in [0, 1]$.
  \item If $z$ has no repeated values, then 
  $\hat{F}_z(\text{Quantile}(z; \tau)) = \frac{\lceil n\tau \rceil}{n}$ for all $\tau \in [0, 1]$.
\end{enumerate}

\paragraph{Solution}
Let $z \in \mathbb{R}^n$ be given and fixed.
Let us inrtoduce a notation $S_{\tau}$
for the set of all $v \in \mathbb{R}$ such that $\hat{F}_z(v) \geqslant \tau$.
Then, by definition, $\text{Quantile}(z; \tau) = \inf S_{\tau}$.
All four statements depend on the following lemma:

\begin{lemma}
If $v \in \mathbb{R}$ and $\tau \in [0, 1]$, then
$\text{Quantile}(z; \tau) \leqslant v$ if and only if $\tau \leqslant \hat{F}_z(v)$.
\end{lemma}

\begin{proof}
If $\tau \leqslant \hat{F}_z(v)$, then $v \in S_{\tau}$ 
and $v \geqslant \inf S_{\tau} = \text{Quantile}(z; \tau)$.

Let us focus on the converse.
Suppose that $\text{Quantile}(z; \tau) \leqslant v$.
For simplicity, put $q = \text{Quantile}(z; \tau)$.
If $\tau = 0$, then $\hat{F}_z(v) \geqslant 0 = \tau$ holds trivially.
Assume from now on that $\tau > 0$.
Then $q$ is finite number.
Choose $\delta > 0$ small enough such that no points 
$z_i$ lies between $q$ and $q + \delta$,
which is always doable because $z$ has only finitely many points.
Since $q$ is the infimum,
choose $s \in S_{\tau}$ such that $q \leqslant s < q + \delta$.
Then, we arrive at 

\begin{equation*}
  \hat{F}_z(v) \geqslant \hat{F}_z(q) = \hat{F}_z(s) \geqslant \tau
\end{equation*}
by applying in the first inequality the monotonicity of $\hat{F}_z$ 
and in the second equality the fact that no points $z_i$ lies between $q$ and $q + \delta$.
The lemma is proved.
\end{proof}

We now prove the four statements.
To prove the first statement, 
let $v \in \mathbb{R}$ be given 
and we have, by the lemma:

\begin{equation*}
  \curly{\tau \in [0,1] \mid \text{Quantile}(z; \tau) \leqslant v}
  = \curly{\tau \in [0,1] \mid \tau \leqslant \hat{F}_z(v)} = [0, \hat{F}_z(v)]
\end{equation*}
and the suprenum of the interval is just $\hat{F}_z(v)$.
The second and the third statements are already included in the lemma.
The fourth statement is trivial given the 
special assumption that $z$ has no repeated values.////
\end{CJK*}
\end{document}